% Ad Atticum 2.23 --- headnote
% ad-atticum-02-23, before 18 October 59 BC, written at Rome

\textit{Cicero to Atticus, written at Rome shortly before 18
October 59 BC, dictated as he walked --- the first letter, he
notes, that Atticus has ever read in another's hand. The voice
needs rest. The letter introduces two of the cipher names that
will mark the late-59 correspondence: \textit{Sampsiceramus}
for Pompey (the dynastic title of the petty king of Emesa whom
Pompey had patronized in the East --- a private joke at
Pompey's expense for setting himself up as an oriental
potentate); and \textit{Ox-Eyed Hera} \textit{[Greek:
Bo\=opis]} for Clodia, the elder of the three Clodian sisters
and wife of Q. Metellus Celer (consul 60 BC), whose
\textit{consanguineus} (kinsman --- and, in the gossip, more
than that) is Clodius himself.}

\textit{The political picture: Pompey is sick of his condition,
longing to be restored to where he fell from, seeking medicine
that Cicero cannot give. The triumviral side, with no enemy in
the field, is growing old before time; the unanimity of public
opinion against them has never been greater. Yet Clodius is
making his threats: denying them to Pompey, parading them to
everyone else. The closing imperative is the most-quoted line
of late-59 correspondence: \textit{si dormis expergiscere, si
stas ingredere, si ingrederis curre, si curris advola} ---
``if you are asleep, wake; if you are standing, step out; if
you are stepping, run; if you are running, fly.'' The matter
is the elections (postponed to October by Bibulus's
obstruction): if Atticus cannot make the elections, he must
at least be at Rome when ``the man'' --- Clodius --- is
declared tribune.}
